\clearpage
\section*{Problem 1}
\addcontentsline{toc}{section}{Problem 1}
\markboth{Problem 1}{Problem 1}

\textbf{Problem Statement (Textbook 4.1.7):} 

Let $A$ have the block form
$$A = \begin{bmatrix} B & C \\ 0 & I \end{bmatrix}$$
in which the blocks are $n \times n$. Prove that if $(B - I)$ is nonsingular, then, for $k \geq 1$,
$$A^k = \begin{bmatrix} B^k & (B^k - I)(B - I)^{-1}C \\ 0 & I \end{bmatrix}$$

\noindent\rule{\textwidth}{0.4pt}

\vspace{1em}

\textbf{Solution:}

To understand the pattern, let's first compute $A^2$ and $A^3$ directly.

\textbf{Computing $A^2$:}

\begin{align*}
A^2 &= \begin{bmatrix} B & C \\ 0 & I \end{bmatrix} \begin{bmatrix} B & C \\ 0 & I \end{bmatrix} \\[6pt]
&= \begin{bmatrix} B^2 & BC + C \\ 0 & I \end{bmatrix}
\end{align*}

\textbf{Computing $A^3$:}

\begin{align*}
A^3 &= A^2 \cdot A = \begin{bmatrix} B^2 & BC + C \\ 0 & I \end{bmatrix} \begin{bmatrix} B & C \\ 0 & I \end{bmatrix} \\[6pt]
&= \begin{bmatrix} B^3 & B^2C + BC + C \\ 0 & I \end{bmatrix}
\end{align*}

Observing the pattern in the upper-right block:
\begin{align*}
k=1: &\quad C \\
k=2: &\quad BC + C \\
k=3: &\quad B^2C + BC + C
\end{align*}

We can factor these as:
\begin{align*}
k=1: &\quad (I)C \\
k=2: &\quad (B + I)C \\
k=3: &\quad (B^2 + B + I)C
\end{align*}

This suggests the general formula: the upper-right block is $(I + B + B^2 + \cdots + B^{k-1})C$, which equals $(B^k - I)(B - I)^{-1}C$ (see Appendix~\ref{appendix:matrix-geometric-series}).

\vspace{1em}

Now we prove this result by induction on $k$.

\textbf{Base Case ($k=1$):}

For $k=1$, we need to verify that 
$$A^1 = \begin{bmatrix} B & C \\ 0 & I \end{bmatrix} = \begin{bmatrix} B^1 & (B^1 - I)(B - I)^{-1}C \\ 0 & I \end{bmatrix}$$

This simplifies to showing that $C = (B - I)(B - I)^{-1}C$, which is trivially true since $(B - I)(B - I)^{-1} = I$.

\textbf{Inductive Step:}

Assume the formula holds for $k = n$:
$$A^n = \begin{bmatrix} B^n & (B^n - I)(B - I)^{-1}C \\ 0 & I \end{bmatrix}$$

We need to prove it holds for $k = n+1$. Computing $A^{n+1} = A^n \cdot A$:

\begin{align*}
A^{n+1} &= \begin{bmatrix} B^n & (B^n - I)(B - I)^{-1}C \\ 0 & I \end{bmatrix} \begin{bmatrix} B & C \\ 0 & I \end{bmatrix} \\[8pt]
&= \begin{bmatrix} B^n \cdot B + (B^n - I)(B - I)^{-1}C \cdot 0 & B^n \cdot C + (B^n - I)(B - I)^{-1}C \cdot I \\ 0 \cdot B + I \cdot 0 & 0 \cdot C + I \cdot I \end{bmatrix} \\[8pt]
&= \begin{bmatrix} B^{n+1} & B^n C + (B^n - I)(B - I)^{-1}C \\ 0 & I \end{bmatrix}
\end{align*}

Now we need to show that the upper-right block equals $(B^{n+1} - I)(B - I)^{-1}C$.

\textbf{Key Identity:}

We claim that $(B^{n+1} - I)(B - I)^{-1}C = B^n C + (B^n - I)(B - I)^{-1}C$.

This identity is equivalent to showing that:
$$(B^k - I)(B - I)^{-1} = I + B + B^2 + \cdots + B^{k-1}$$

The proof of this geometric series identity for matrices is provided in Appendix~\ref{appendix:matrix-geometric-series}.

Thus, we have shown that:
$$A^{n+1} = \begin{bmatrix} B^{n+1} & (B^{n+1} - I)(B - I)^{-1}C \\ 0 & I \end{bmatrix}$$

By the principle of mathematical induction, the formula holds for all $k \geq 1$.

\vspace{2em}
\noindent\rule{\textwidth}{0.4pt}
